\documentclass[11pt]{amsart}

\usepackage[T1]{fontenc}
\usepackage{lmodern}
\usepackage{microtype}
\usepackage{mathtools,amssymb,amsthm}
\usepackage[hidelinks]{hyperref}

\hypersetup{
  pdftitle={Under a simplicity hypothesis, the order of a [3,1;4]-mixed cage is twelve}
}

\newtheorem{theorem}{Theorem}
\newtheorem{lemma}[theorem]{Lemma}
\newtheorem{corollary}[theorem]{Corollary}
\theoremstyle{remark}
\newtheorem*{remark}{Remark}

\newcommand{\mn}[1]{n[#1]}

\title[{Order twelve under a simplicity hypothesis}]{Under a simplicity
hypothesis, the order of a $[3,1;4]$-mixed cage is twelve}

\date{}

\subjclass[2020]{05C35, 05C20, 05C38}
\keywords{mixed graph, mixed cage, girth, directed cage, circulant}

\begin{document}

\begin{abstract}
Araujo-Pardo, Hern\'andez-Cruz and Montellano-Ballesteros proved
$\mn{z,1;4}\le 3(z+1)$ for odd $z$ and $\mn{z,1;4}\le 3z+2$ for even $z$, with
equality for $z\in\{1,2\}$, and proposed as a first problem the search for a
matching lower bound. Under the standing hypothesis that mixed graphs are
\emph{simple}, we compute the single case $z=3$: $\mn{3,1;4}=12$. The upper half
is their own circulant on $\mathbb Z_{12}$; the lower half is two lines from an
elementary digraph lemma --- a loopless, digon-free, directed-triangle-free
digraph on $m$ vertices with minimum out-degree $\delta\ge1$ has $m\ge
2\delta+2$ --- applied to the digraph induced on the out-neighbourhood of one
vertex. The same argument gives $\mn{z,1;4}\ge\lceil(5z+6)/2\rceil$ for all
$z\ge2$. The simplicity reading is a hypothesis, not a theorem about the source,
whose text contains both a sentence allowing multiple edges and arcs and a
definition requiring a simple graph; we exhibit an order-$6$ multigraph object
for which the value $12$ fails, so the multigraph variant is not decided here.
\end{abstract}

\maketitle

\section{The problem}

A \emph{mixed graph} has both edges and arcs. Following Araujo-Pardo,
Hern\'andez-Cruz and Montellano-Ballesteros \cite{AHM}, a mixed graph is
\emph{$[z,r]$-regular} if every vertex is the tail of exactly $z$ arcs, the head
of exactly $z$ arcs, and incident with exactly $r$ edges; a \emph{walk}
traverses an edge in either direction and an arc only forwards; the \emph{girth}
is the length of a shortest cycle in that sense; and $\mn{z,r;g}$ denotes the
least order of a $[z,r]$-regular mixed graph of girth $g$, a \emph{$[z,r;g]$-mixed
cage}. Throughout, mixed graphs are \emph{simple}: at most one arc joins a given
ordered pair of vertices, and no arc is parallel to an edge. This convention is
a hypothesis, and it is load-bearing; we return to it below. Since $r=1$ means
the edges form a perfect matching, the order of a $[z,1;g]$-mixed graph is even.

The target statement is Theorem~3.3 of \cite{AHM}, verbatim:

\begin{quote}
\textbf{Theorem.} Let $z\ge1$ be an integer, then we have that
$\mn{z,1;4}\le 3(z+1)$ if $z$ is odd, $3z+2$ otherwise. Moreover, the equality
holds for $z\in\{1,2\}$.
\end{quote}

The equality clause covers $z\in\{1,2\}$; the case treated here is $z=3$, where
$3(z+1)=12$. A matching lower bound, in general and not specifically at $z=3$,
is asked for in the closing paragraph of Section~4 of \cite{AHM}:

\begin{quote}
As a first problem, we propose to find a lower bound for the order of a
$[z,1;4]$-mixed cage. This, together with the upper bound given in Theorem~3.3,
could result in an exact value for the order of this family of mixed cages.
\end{quote}

\medskip
\noindent\textbf{The simplicity convention.} Two clauses of \cite{AHM} conflict:
its opening says ``We allow multiple edges and arcs'', while its definition of a
mixed regular graph asks for ``a simple and finite graph $G$''. We do not resolve
which reading the problem quoted above intends, and we make no claim to settle
or extend that problem. Instead we prove a simple-graph statement, with
simplicity as an explicit hypothesis; the remark in Section~3 shows that the
value below is false under the multigraph reading.

\begin{theorem}\label{thm:main}
For simple mixed graphs, $\mn{3,1;4}=12$.
\end{theorem}

Consequently, for simple mixed graphs, the odd-$z$ bound $3(z+1)$ of
Theorem~3.3 of \cite{AHM} is tight at $z=3$. Whether this is the same assertion
as an extension of that theorem's equality clause depends on the unresolved
convention discussed above.

\section{The witness}

The inequality $\mn{3,1;4}\le 12$ is Lemma~3.1(1) of \cite{AHM} at $n=3$: the
mixed circulant $C$ on $\mathbb Z_{12}$ with

\begin{center}
arcs $i\to i+1$, $i\to i+2$, $i\to i+7$ \qquad and \qquad edges $\{i,i+6\}$,
\end{center}

all indices modulo $12$. In full:

\begin{center}\ttfamily\small
\begin{tabular}{l}
a 0 1\ \ a 0 2\ \ a 0 7\ \ a 1 2\ \ a 1 3\ \ a 1 8\ \ a 2 3\ \ a 2 4\ \ a 2 9\\
a 3 4\ \ a 3 5\ \ a 3 10\ \ a 4 5\ \ a 4 6\ \ a 4 11\ \ a 5 0\ \ a 5 6\ \ a 5 7\\
a 6 1\ \ a 6 7\ \ a 6 8\ \ a 7 2\ \ a 7 8\ \ a 7 9\ \ a 8 3\ \ a 8 9\ \ a 8 10\\
a 9 4\ \ a 9 10\ \ a 9 11\ \ a 10 0\ \ a 10 5\ \ a 10 11\ \ a 11 0\ \ a 11 1\ \ a 11 6\\[2pt]
e 0 6\ \ e 1 7\ \ e 2 8\ \ e 3 9\ \ e 4 10\ \ e 5 11
\end{tabular}
\end{center}

\noindent(\texttt{a u v} is the arc $u\to v$; \texttt{e u v} is the edge
$\{u,v\}$; $36$ arcs and $6$ edges on $12$ vertices.)

Out-degree and in-degree are $3$, the in-differences
being $12-1=11$, $12-2=10$, $12-7=5$. The edges form a perfect matching, so
$r=1$. There is no digon, since none of $11,10,5$ lies in $\{1,2,7\}$; no arc is
parallel to an edge, since $6\notin\{1,2,7\}$. There is no directed triangle,
because the ten three-term sums from $\{1,2,7\}$ with repetition are
\[
 3,\;4,\;9,\;5,\;10,\;3,\;6,\;11,\;4,\;9 \pmod{12},
\]
none of them $0$; and no mixed triangle (an edge $\{x,y\}$ together with arcs
$x\to w\to y$), because the six two-term sums
\[
 2,\;3,\;8,\;4,\;9,\;2 \pmod{12}
\]
avoid the edge difference $6$. Finally $1+2+2+7=12$, so $0\to1\to3\to5\to0$ is a
directed $4$-cycle. Hence the girth is exactly $4$ and $C$ is a
$[3,1;4]$-mixed graph of order $12$.

\section{The lower bound}

Everything rests on one elementary lemma about digraphs.

\begin{lemma}\label{lem:L}
Let $H$ be a digraph on $m$ vertices with no loop, no digon, no directed
triangle, and minimum out-degree $\delta\ge1$. Then $m\ge 2\delta+2$.
\end{lemma}

\begin{proof}
The number of arcs is at least $m\delta$ and equals the sum of the in-degrees,
so some vertex $a$ has in-degree at least $\delta$. Put $Q=N^-(a)$ and
$P=N^+(a)$; then $|Q|,|P|\ge\delta\ge1$, and $P\cap Q=\varnothing$ because a
vertex in both would give a digon at $a$, and $a\notin P\cup Q$ because there is
no loop. Pick $p\in P$. Then $N^+(p)$ omits $p$ (no loop), omits $a$ (else
$a\to p\to a$ is a digon), and omits every $q\in Q$ (else $a\to p\to q\to a$ is
a directed triangle on three distinct vertices). Hence
\[
 \delta\;\le\;|N^+(p)|\;\le\;m-2-|Q|\;\le\;m-2-\delta . \qedhere
\]
\end{proof}

Lemma~\ref{lem:L} is tight at $\delta=1$: the directed $4$-cycle has $m=4=2\cdot
1+2$. That is the case used below. At $m=3$ the lemma says there is no such
digraph with $\delta\ge1$, and one sees it directly: follow $a\to b\to c$; then
$c\to a$ is a directed triangle and $c\to b$ a digon, so $c$ has nowhere to
point.

\begin{theorem}\label{thm:lower}
For every $z\ge2$, $\mn{z,1;4}\ge\lceil(5z+6)/2\rceil$, and hence
$\mn{z,1;4}\ge 2\lceil(5z+6)/4\rceil$ because $r=1$ forces the order to be even.
\end{theorem}

\begin{proof}
Let $(D,M)$ be a $[z,1;4]$-mixed graph of order $n$, with $M$ the perfect
matching of edges. Fix a vertex $u$, let $v=M(u)$, $A=N^+(u)$, $B=N^-(u)$ and
$S=V\setminus(\{u,v\}\cup A\cup B)$. These four sets are pairwise disjoint:
$A\cap B=\varnothing$ since a common vertex would give a digon, $u\notin A\cup B$
since there is no loop, and $v\notin A\cup B$ since an arc parallel to the edge
$\{u,v\}$ is a cycle of length $2$. As $|A|=|B|=z$, we get
$|S|=n-2z-2$.

Let $a\in A$. Then $N^+(a)$ omits $u$ (a digon with $u\to a$), omits $v$
(otherwise $u\to a\to v$ together with the edge $\{v,u\}$ is a cycle of length
$3$), and omits every $b\in B$ (otherwise $u\to a\to b\to u$ is a directed
triangle). So $N^+(a)\subseteq (A\setminus\{a\})\cup S$, whence
$|N^+(a)\cap A|\ge z-|S|$.

Let $H$ be the digraph induced on $A$. It has $m=|A|=z$ vertices, and it
inherits from $D$ the absence of loops, digons and directed triangles. If
$z-|S|\le0$ then $|S|\ge z$ and $n=2z+2+|S|\ge 3z+2$. Otherwise $H$ has
minimum out-degree at least $z-|S|\ge1$, so Lemma~\ref{lem:L} gives
$z\ge 2(z-|S|)+2$, i.e. $|S|\ge\lceil (z+2)/2\rceil$ and therefore
\[
 n=2z+2+|S|\;\ge\;2z+2+\Bigl\lceil\tfrac{z+2}{2}\Bigr\rceil
 \;=\;\Bigl\lceil\tfrac{5z+6}{2}\Bigr\rceil .
\]
For $z\ge2$ the second alternative is the smaller of the two, so it is the
bound. The parity refinement is immediate.
\end{proof}

\begin{proof}[Proof of Theorem~\ref{thm:main}]
Suppose a $[3,1;4]$-mixed graph has order $10$. Then, in the notation of the
previous proof, $|S|=10-2\cdot3-2=2$, so the digraph induced on $A=N^+(u)$ has
$3$ vertices, no loop, no digon, no directed triangle and minimum out-degree at
least $3-2=1$. Lemma~\ref{lem:L} demands $3\ge 4$, a contradiction. Orders
below $10$ are excluded by the same one-vertex count: $n\le 6$ is impossible
because $\{u\}\cup\{v\}\cup A\cup B$ already has $2+3+3=8$ elements, and $n=8$
forces $|S|=0$, so $H$ would have minimum out-degree $3$ on $3$ vertices and
Lemma~\ref{lem:L} demands $3\ge8$. As the order is even, $\mn{3,1;4}\ge12$;
equivalently, Theorem~\ref{thm:lower} at $z=3$ gives
$\lceil 21/2\rceil=11$ and parity lifts it to $12$. The circulant of
Section~2 has order $12$. Hence $\mn{3,1;4}=12$.
\end{proof}

At $z=2$, Theorem~\ref{thm:lower} returns $\lceil16/2\rceil=8$, the published
value $\mn{2,1;4}=8$ of \cite{AHM}.

\begin{remark}[the simplicity hypothesis is load-bearing]
Without it the value drops to at most $6$. On $\mathbb Z_6$ take the arc
$i\to i+1$ with multiplicity three and the edges $\{i,i+3\}$: every vertex is the
tail of three arcs, the head of three arcs and incident with one edge; parallel
arcs point the same way, so they create no cycle of length $2$; $3\notin\{1\}$,
so no arc is parallel to an edge; $1+1+1=3\ne0\bmod 6$, so there is no directed
triangle; $1+1=2\ne3$, so there is no mixed triangle; and $0\to1\to2\to3$ with
the edge $\{3,0\}$ is a cycle of length $4$. The multigraph variant is not
decided here.
\end{remark}

\section{Prior work}

Delete the perfect matching from a $[z,1;4]$-mixed graph and there remains a
digraph of minimum out-degree $z$ with no digon and no directed triangle; counts
of that kind are standard: the introduction of \cite{AHM} records that Behzad,
Chartrand and Wall \cite{BCW} proved the analogous lower bound
$\lceil(5z+4)/2\rceil$ on the order of a $z$-regular digraph of girth $4$. We
claim no novelty for Lemma~\ref{lem:L} or for the count in
Theorem~\ref{thm:lower}.

We do not claim that the circulant of Section~2 is the unique simple
$[3,1;4]$-mixed cage, and we say nothing about $g\ne4$ or $r\ne1$.

\section{Verification}

The mathematics needs no machine: the order-$12$ witness is settled by the
difference computations of Section~2, and Lemma~\ref{lem:L} together with the
count in the proof of Theorem~\ref{thm:main} excludes every smaller order in two
lines. The accompanying program \texttt{verify.py} re-derives that in exact
integer arithmetic from the link list printed in Section~2, confirms
Lemma~\ref{lem:L} by exhaustive enumeration of all loopless, digon-free,
directed-triangle-free digraphs on at most five vertices, and rechecks the objects
exhibited in Sections~2 and~3.
Its transcript \texttt{verify.output.txt} records $66$ passing checks in nine
parts; parts~1--5 and~7 are the ones bearing on the statements above, and
parts~6, 8 and~9 are arithmetic on constants quoted from the literature and
comparisons no longer discussed here. The transcript's scope note names the
quoted constants, and records that nothing there decides the multigraph variant
beyond exhibiting the order-$6$ object, and that nothing there claims the
order-$12$ witness is unique.

\begin{thebibliography}{9}

\bibitem{AHM}
G.~Araujo-Pardo, C.~Hern\'andez-Cruz, and J.~J. Montellano-Ballesteros,
\emph{Mixed cages},
Graphs Combin. \textbf{35} (2019), no.~4, 989--999.
\href{https://doi.org/10.1007/s00373-019-02050-1}{doi:10.1007/s00373-019-02050-1};
\href{https://arxiv.org/abs/1702.07255}{arXiv:1702.07255}.
Quotations above are from the preprint.

\bibitem{BCW}
M.~Behzad, G.~Chartrand, and C.~E. Wall,
\emph{On minimal regular digraphs with given girth},
Fund. Math. \textbf{69} (1970), 227--231.

\end{thebibliography}

\end{document}
